\documentclass{article}

\usepackage{amsmath, amssymb, geometry}
\usepackage{amsthm}
\usepackage{hyperref}

\begin{document}

\begin{flushleft}
    \textbf{\Large Tiago Lacerda} \\[4pt]
    Álgebra Linear Computacional \\[2pt]
    Lista de Exercícios 03 \\[2pt]
    2026.1
\end{flushleft}

4. Seja $P$ a matriz de permutação de linhas da decomposição LU com
pivoteamento total. Por construção, $P$ é uma matriz em que cada linha e cada
coluna contém exatamente um elemento igual a $1$, sendo todos os demais
elementos iguais a $0$.

Assim, toda linha e toda coluna de $P$ tem norma $1$, e o produto interno entre
linhas ou colunas distintas é $0$. Considere:

\[
    (P^{T}P)_{ij} = \sum_{k} P_{ki}P_{kj}.
\]

Se $i=j$, a soma resulta em $1$; caso contrário, resulta em $0$. Portanto,

\[
    P^{T}P = I,
\]

ou seja, $P^{-1} = P^{T}$, atendendo à definição de matriz ortogonal.

O mesmo argumento é válido para a matriz de permutação de colunas $Q$.

\bigbreak

5. De maneira geral, imagens naturais não apresentam grandes variações entre pixeis vizinhos, exceto em contornos e bordas. Assim, bases como as de Fourier e as wavelets de Haar são capazes de representar esses dados com uma quantidade reduzida de coeficientes sem grande perda de fidelidade. O nível de compressão pode então ser facilmente controlado pelo número de coeficientes mantidos após a transformação. Além disso, por serem bases ortogonais, suas transformações inversas são dadas pelas transpostas das matrizes de transformação, tornando a reconstrução das imagens direta, sem cálculo da inversa de uma matriz genérica.

\bigbreak

9, 10. https://github.com/TiagoLacerda/alg-lin-comp-lista-3






\end{document}
